% 谷神星（综述）
% license CCBYSA3
% type Wiki

本文根据 CC-BY-SA 协议转载翻译自维基百科\href{https://en.wikipedia.org/wiki/Ceres_(dwarf_planet)}{相关文章}。

\begin{figure}[ht]
\centering
\includegraphics[width=6cm]{./figures/ed4afc93bd247d5b.png}
\caption{Dawn 探测器于 2015 年 5 月拍摄的 Ceres 图像。两个明亮斑点分别为 Haulani 撞击坑（右）和 Oxo 撞击坑坑底（左）。} \label{fig_Ceres_1}
\end{figure}
Ceres（小行星编号：1 Ceres）是一颗位于 Mars 与 Jupiter 轨道之间主小行星带中的矮行星。它是人类发现的首颗小行星，于 1801 年 1 月 1 日由 Giuseppe Piazzi 在 Sicily 的 Palermo Astronomical Observatory 发现，并在当时被宣布为一颗新行星。此后，Ceres 被重新分类为小行星，近年又被认定为矮行星——它是唯一一颗轨道未超出 Neptune 之外的矮行星，也是所有不具备卫星的矮行星中最大的一颗。

Ceres 的直径约为 Moon 的四分之一。由于其体积很小，即使在最亮时，除非在极其黑暗的天空条件下，也无法用肉眼看到。其视星等范围为 6.7 至 9.3，并在每隔 15 至 16 个月发生一次会合（即靠近 Earth）时达到峰值。因此，即使使用最强大的望远镜，其表面特征也只能勉强识别，在 NASA Dawn 无人探测器于 2015 年抵达并进入 Ceres 轨道任务之前，人们对其所知甚少。

Dawn 的观测发现，Ceres 的表面由水、冰以及碳酸盐和黏土等含水矿物混合组成。重力数据表明，Ceres 的内部可能呈部分分化结构，拥有一个泥状（冰—岩）地幔/核心，以及一层密度较低但更坚固的地壳（其体积最多约 30\% 为冰）。虽然 Ceres 可能缺乏内部液态水海洋，但盐水仍在外地幔中流动并抵达表面，从而大约每五千万年形成一次如 Ahuna Mons 这样的低温火山。这使 Ceres 成为距离 Sun 最近的已知低温火山活跃天体。Ceres 具有极其稀薄且短暂的水蒸气大气，由其表面局部区域喷发形成。
\subsection{历史}  
\subsubsection{发现}  
在 18 世纪日心说被接受到 1846 年 Neptune 被发现之间的多年里，一些天文学家认为数学定律预测在 Mars 和 Jupiter 轨道之间存在一颗隐藏或失踪的行星。1596 年，理论天文学家 Johannes Kepler 认为，只有在添加两颗行星——一颗位于 Jupiter 与 Mars 之间，另一颗位于 Venus 与 Mercury 之间——时，行星轨道之间的比例才会符合“上帝的设计”\(^{[19]}\)。其他理论家，例如 Immanuel Kant，思考该空隙是否由 Jupiter 的引力造成；1761 年，天文学家兼数学家 Johann Heinrich Lambert 问道：“谁又能知道是否已有行星从 Mars 与 Jupiter 之间的广阔空间中消失？天体是否如同 Earth 一样，强者磨损弱者，Jupiter 与 Saturn 是否注定永远掠夺它们？”\(^{[19]}\)

1772 年，德国天文学家 Johann Elert Bode 引用 Johann Daniel Titius 的研究，发表了一条后来被称为 Titius–Bode 定律的公式，该公式似乎能预测已知行星的轨道，但在 Mars 与 Jupiter 之间出现无法解释的空隙\(^{[19][20]}\)。该公式预测 Sun 附近约 2.8 AU（约 4.2 亿 km）处应存在另一颗行星\(^{[20]}\)。随着 William Herschel 于 1781 年在预测距离附近发现 Saturn 之外的 Uranus，这一定律获得更多信任\(^{[19]}\)。1800 年，由德国天文期刊《Monthly Correspondence》主编 Franz Xaver von Zach 领导的一组人向 24 名经验丰富的天文学家发出请求，将他们称为“天上警察”\(^{[20]}\)，希望他们共同努力，开始系统搜索这颗预期行星\(^{[20]}\)。尽管他们未发现 Ceres，但后来发现了小行星 Pallas、Juno 和 Vesta\(^{[20]}\)。

受邀参与该搜索计划的天文学家之一是 Giuseppe Piazzi，他是 Sicily Palermo 学院的一名天主教神父。在收到邀请函之前，Piazzi 于 1801 年 1 月 1 日发现了 Ceres\(^{[21]}\)。他当时在寻找“la Caille 氏黄道星表中的第 87 颗恒星”，却发现“前面还有另一颗”\(^{[19]}\)。Piazzi 并非找到了一颗恒星，而是一颗运动中的类星体，他起初认为那是一颗彗星\(^{[22]}\)。Piazzi 共观测 Ceres 24 次，最后一次观测发生在 1801 年 2 月 11 日，后因生病中断。他于 1801 年 1 月 24 日通过信件向两位天文学家宣布了他的发现：一位是来自 Milan 的同胞 Barnaba Oriani，另一位是 Berlin 的 Bode\(^{[23]}\)。他将其报告为彗星，但表示“由于它的运动如此缓慢且近乎均匀，我多次想到它可能比彗星更重要”\(^{[19]}\)。4 月，Piazzi 将完整观测结果送给 Oriani、Bode 和法国天文学家 Jérôme Lalande。这些信息发表于 1801 年 9 月的《Monatliche Correspondenz》\(^{[22]}\)。

当时，Ceres 的视位置已经发生改变（主要由于 Earth 绕 Sun 运动），且过于靠近 Sun 的眩光，其他天文学家无法确认 Piazzi 的观测。到了年底，Ceres 应该能再次被观测到，但由于已经过去很长时间，其位置难以预测。为重新定位 Ceres，24 岁的数学家 Carl Friedrich Gauss 开发了一套高效的轨道确定方法\(^{[22]}\)。他在几周内预测出 Ceres 的轨迹，并将结果发送给 von Zach。1801 年 12 月 31 日，von Zach 与同为“天上警察”的 Heinrich W. M. Olbers 在预测位置附近再次发现 Ceres，并继续记录其位置\(^{[22]}\)。Ceres 与 Sun 相距 2.8 AU，几乎完美符合 Titius–Bode 定律；但当 1846 年 Neptune 被发现时，其实际位置比预测近 8 AU，多数天文学家认为该定律仅是巧合\(^{[24]}\)。

早期观测者只能大致估算 Ceres 的大小。Herschel 在 1802 年低估其直径为 260 km（160 mi）；1811 年，德国天文学家 Johann Hieronymus Schröter 将其高估为 2,613 km（1,624 mi）\(^{[25]}\)。20 世纪 70 年代，红外光度测量使其反照率测定更为精确，Ceres 的直径得出与真实值 939 km（583 mi）相差不到 10\% 的结果\(^{[25]}\)。
\subsubsection{名称与符号}  
Piazzi 为其发现所提出的名称是 Ceres Ferdinandea：“Ceres”取自罗马农业女神，其神庙最早建立于 Sicily；“Ferdinandea”则是为了纪念 Piazzi 的君主与资助者——Sicily 的 Ferdinand III 国王\(^{[22]}\)。后一个名称未被其他国家接受，因而被取消。在 von Zach 于 1801 年 12 月重新观测到 Ceres 之前，von Zach 将该行星称为 Hera，而 Bode 则称其为 Juno。尽管 Piazzi 表示反对，这些名称在德国一度通行，直到该天体的存在得到确认。确认之后，天文学家最终采用了 Piazzi 所命名的名称\(^{[26]}\)。

Ceres 的形容词形式有 Cererian\(^{[27][28]}\) 和 Cererean\(^{[29]}\)，两者读音均为 /sɪˈrɪəriən/\(^{[30][31]}\)。1803 年发现的一种稀土元素 cerium（铈）即以 Ceres 命名\(^{[32][c]}\)。

Ceres 的旧天文符号今天仍在占星术中使用，为一柄镰刀，⟨⚳⟩\(^{[22][34]}\)。镰刀是女神 Ceres 的经典象征之一，该符号显然分别由 von Zach 与 Bode 于 1802 年独立提出\(^{[35]}\)。它的形状与 Venus 的符号相似（⟨♀⟩，即一个圆圈下加小十字），但在圆圈处有一处断开。其表现形式存在多种细微变体，包括一种反向图形，以字母 “C”（即 Ceres 的首字母）配以一个加号进行排版。编号圆盘形式的小行星通用符号 ① 于 1867 年引入，并迅速成为惯例\(^{[22][36]}\)。
\subsubsection{分类}
\begin{figure}[ht]
\centering
\includegraphics[width=6cm]{./figures/f71c683b033e055f.png}
\caption{按比例显示的 Ceres（左下）、Moon 与 Earth} \label{fig_Ceres_2}
\end{figure}
\begin{figure}[ht]
\centering
\includegraphics[width=6cm]{./figures/7ea9a2828e334e64.png}
\caption{小行星带中四颗最大小行星的平均直径对比（左侧为矮行星 Ceres）} \label{fig_Ceres_3}
\end{figure}
Ceres 的分类经历过多次变化，并引发过不少争议。Bode 认为 Ceres 是他所提出的位于 Mars 与 Jupiter 之间“失踪行星”\(^{[19]}\)。Ceres 被赋予了行星符号，并在天文学书籍与星表中（与 Pallas、Juno 和 Vesta 一起）作为行星列出，持续了半个多世纪\(^{[37]}\)。

随着更多天体在 Ceres 附近被发现，天文学家开始怀疑它代表了一类新的天体\(^{[19]}\)。当 Pallas 于 1802 年被发现时，Herschel 引入了 “asteroid”（“类星状”）这一术语用于这些天体\(^{[37]}\)，写道：“它们与小星星极为相似，即便使用非常好的望远镜也几乎无法区别”\(^{[38]}\)。1852 年，Johann Franz Encke 在《Berliner Astronomisches Jahrbuch》中宣布，为这些新天体授予传统的行星符号体系过于繁琐，因而引入依照发现顺序在名称前添加数字的新方法。编号体系最初从第五颗小行星 5 Astraea 作为 1 开始，但在 1867 年，Ceres 被纳入该体系，命名为 1 Ceres\(^{[37]}\)。

到 1860 年代，天文学家普遍接受了行星与诸如 Ceres 之类小行星之间存在根本区别，尽管当时 “planet（行星）” 一词尚未被精确定义\(^{[37]}\)。在 1950 年代，科学家普遍停止将大多数小行星视为行星，但 Ceres 有时仍保留这一地位，因为其在地球物理复杂性方面具有类行星特征\(^{[39]}\)。随后在 2006 年，围绕 Pluto 的争议引发了对“行星”定义的呼声，并提出重新分类 Ceres，甚至可能恢复其作为行星的普遍地位\(^{[40]}\)。一项提交给国际天文学联合会（IAU，负责天文命名与分类的全球组织）的提案，将行星定义为：“（a）具有足够质量，其自身引力能克服刚体力使其呈流体静力平衡（近似球形）形状，并且（b）围绕恒星运行，且既非恒星亦非行星的卫星”\(^{[41]}\)。若该决议被采纳，则 Ceres 将成为距 Sun 第五颗行星\(^{[42]}\)，但 2006 年 8 月 24 日大会增加了额外要求，即行星必须“清除其轨道附近的区域”。Ceres 并非行星，因为它无法主导其轨道，与小行星带中成千上万的其他小行星共享轨道，其质量仅占小行星带总质量的约 40\%\(^{[43]}\)。那些满足最初定义但不满足第二条要求的天体（如 Ceres），被分类为矮行星\(^{[44]}\)。行星地质学家仍常常无视这一定义，并依然将 Ceres 视为行星\(^{[45]}\)。

Ceres 被指定为矮行星与小行星。NASA 一篇网页写道，小行星带中第二大天体 Vesta 是最大的小行星\(^{[46]}\)。IAU 在这一问题上态度模糊\(^{[47][48][\text{failed verification}]}\)，但其负责编目相关天体的 Minor Planet Center 指出，矮行星可能具有双重指定\(^{[49]}\)，而 IAU / USGS / NASA 联合地名索引将 Ceres 同时归类为小行星与矮行星\(^{[50]}\)。
\begin{figure}[ht]
\centering
\includegraphics[width=6cm]{./figures/df7ba81daeff6318.png}
\caption{} \label{fig_Ceres_4}
\end{figure}
\subsection{轨道}
\begin{figure}[ht]
\centering
\includegraphics[width=8cm]{./figures/0f8d0cf9a4498022.png}
\caption{Ceres（红色，倾斜）的轨道以及 Jupiter 和内行星（白色与灰色）的轨道。上方示意图从俯视角展示了 Ceres 的轨道。下方示意图为侧视图，显示 Ceres 轨道相对于黄道面的倾角。浅色表示位于黄道面之上；深色表示位于黄道面之下。} \label{fig_Ceres_5}
\end{figure}
Ceres 在 Mars 与 Jupiter 之间运行，位于小行星带中部，其轨道周期为 4.6 个 Earth 年\(^{[2]}\)。与其他行星和矮行星相比，Ceres 的轨道相对 Earth 有适度倾斜；其轨道倾角（\(i\)）为 \(10.6^\circ\)，相比之下 Mercury 为 \(7^\circ\)，Pluto 为 \(17^\circ\)。其轨道也略微拉长，偏心率（\(e\)）为 0.08，而 Mars 为 0.09\(^{[2]}\)。

Ceres 不属于任何小行星族，可能是因为其内部拥有较高比例的冰，具有相同成分的较小天体在太阳系年龄范围内会由于升华而消失殆尽\(^{[51]}\)。曾有人认为 Ceres 属于 Gefion 小行星族\(^{[52]}\)，该族成员具有相似的本征轨道要素，暗示它们可能源于过去一次小行星碰撞。然而，后来发现 Ceres 与 Gefion 家族在成分上不同\(^{[52]}\)，似乎属于“闯入者”，即轨道参数相似但并无共同起源\(^{[53]}\)。
\subsubsection{共振}  
由于小行星带天体质量较小且彼此距离较大，它们之间很少出现引力共振\(^{[54]}\)。尽管如此，Ceres 能够将其他小行星暂时捕获进入 1:1 共振轨道（使其成为暂时的特洛伊天体），持续时间从几十万年至两百多万年不等。目前已识别出 50 个这样的天体\(^{[55]}\)。Ceres 与 Pallas 之间接近 1:1 平均运动轨道共振（其本征轨道周期相差 0.2\%），但不足以在天文学时间尺度上产生显著影响\(^{[56]}\)。
\subsection{自转与轴倾角}
Ceres 的自转周期（即 Ceres 一日）为 9 小时 4 分钟\(^{[10]}\)；其赤道上的小型撞击坑 Kait 被选定为本初子午线\(^{[57]}\)。Ceres 的自转轴倾角为 \(4^\circ\)\(^{[10]}\)，这一数值足够小，使其极区包含永远处于阴影中的撞击坑，这些撞击坑预计会充当冷阱，随时间积累水冰，类似于 Moon 和 Mercury 上发生的情况。预计约 0.14\% 从表面释放的水分子最终会被捕获在这些冷阱中，它们平均“跳跃”三次才会逃逸或被捕获\(^{[10]}\)。

首个进入 Ceres 轨道的探测器 Dawn 确定，Ceres 的北极轴指向赤经 19 h 25 m 40.3 s（291.418°）、赤纬 +66° 45′ 50″（距离 Delta Draconis 约 1.5°），这意味着自转轴倾角为 \(4^\circ\)。这表明 Ceres 目前几乎不存在随纬度变化的季节性阳光差异\(^{[58]}\)。在过去三百万年中，Jupiter 与 Saturn 的引力影响引发了 Ceres 自转轴倾角的周期性变化，范围从 \(2^\circ\) 到 \(20^\circ\)，这意味着过去曾存在季节性日照变化，最近一次出现季节活动大约发生在 14,000 年前。在自转轴最大倾角时期仍能保持阴影的那些撞击坑最有可能在太阳系年龄尺度内保留由喷发或彗星撞击产生的水冰\(^{[59]}\)。
\subsection{地质}
\begin{figure}[ht]
\centering
\includegraphics[width=8cm]{./figures/23fb842f59355847.png}
\caption{谷神星在内太阳系 行星质量物体,按它们从太阳向外的轨道顺序排列 (从左起:水银,维纳斯,地球,the月亮,火星和谷神星)} \label{fig_Ceres_6}
\end{figure}
\begin{figure}[ht]
\centering
\includegraphics[width=8cm]{./figures/18dd605eb8e54f28.png}
\caption{Ceres渲染与搅拌机,从色基图和高度图从NASA和USGS。} \label{fig_Ceres_7}
\end{figure}
谷神星是小行星主带中最大的天体$^{[16]}$。它被分类为 C 型或碳质小行星$^{[16]}$，并且由于存在粘土矿物，还被归类为 G 型小行星$^{[60]}$。它的成分与碳质球粒陨石类似，但并不完全相同$^{[61]}$。它是一个扁球体，其赤道直径比极地直径大 8\%$^{[2]}$。“黎明号”探测器的测量发现，谷神星的平均直径为 939.4 公里（583.7 英里）$^{[2]}$，质量为 $9.38\times10^{20}$ 千克$^{[62]}$。这使得谷神星的密度为 2.16 克/立方厘米$^{[2]}$，表明其约四分之一的质量为水冰$^{[63]}$。

谷神星占据小行星带估计总质量 $(2394\pm5)\times10^{18}$ 千克的 40\%，并且其质量是下一颗小行星——灶神星的 3 又 1/2 倍，但仅为月球质量的 1/78，其表面重力为地球的 1/35（约为月球的 1/6）。它非常接近处于流体静力平衡状态，但仍有一些偏离平衡形状的现象尚未得到解释$^{[64]}$。谷神星是唯一公认轨道周期短于海王星的矮行星$^{[63]}$。建模显示，谷神星的岩石物质可能发生了部分分异，并且可能拥有一个小型核心$^{[65]}$,$^{[66]}$，但数据也与一个由水合硅酸盐构成的地幔且没有核心的结构相一致$^{[64]}$。由于“黎明号”没有配备磁强计，因此目前尚不清楚谷神星是否具有磁场；普遍认为其不具备磁场$^{[67]}$,$^{[68]}$。谷神星内部的分异可能与其缺乏天然卫星有关，因为主带小行星的卫星通常被认为是由碰撞破坏形成，从而产生未分化的碎石堆结构$^{[69]}$。

\subsubsection{表面}
\textbf{成分}
谷神星的表面成分在全球尺度上是均一的，并且富含碳酸盐和经水改变的含铵层状硅酸盐$^{[64]}$，尽管风化层中的水冰含量在极地纬度约为 10\%，而在赤道地区则要干燥得多，甚至几乎不含冰$^{[64]}$。

利用哈勃空间望远镜的研究显示，在谷神星表面存在石墨、硫和二氧化硫。石墨显然是谷神星较古老表面经历空间风化的结果；而后两者在谷神星条件下是挥发性的，理论上应当迅速逃逸或沉积于冷阱中，因此它们显然与相对近期的地质活动有关$^{[70]}$。

在厄努特陨石坑检测到有机化合物$^{[71]}$，并且另有至少 11 个区域被认为可能存在有机化合物$^{[72]}$。该行星近表层大部分区域富含碳，质量分数约为 20\%$^{[73]}$。其碳含量比在地球上分析的碳质球粒陨石高出五倍以上$^{[73]}$。表面碳显示出与岩石—水相互作用产物（如粘土）混合的证据$^{[73]}$。这种化学性质表明谷神星形成于寒冷环境中，可能位于木星轨道之外，并且在水的存在下由超高含碳物质聚合而成，从而可能提供有利于有机化学的条件$^{[73]}$。
\begin{figure}[ht]
\centering
\includegraphics[width=8cm]{./figures/97253d100f841e67.png}
\caption{黑白的谷神星摄影地图，以 180° 经线为中心，并标注官方命名（2017 年 9 月）。} \label{fig_Ceres_8}
\end{figure}
\begin{figure}[ht]
\centering
\includegraphics[width=8cm]{./figures/c0d68798ceb10be6.png}
\caption{谷神星，极地区域（2015 年 11 月）：北极（左）；南极（右）。南极处于阴影中。“Ysolo 山”随后被更名为 “Yamor 山”$^{[74]}$。} \label{fig_Ceres_10}
\end{figure}
\textbf{陨石坑}
\begin{figure}[ht]
\centering
\includegraphics[width=8cm]{./figures/8310f7c97e44cd51.png}
\caption{谷神星的地形图。最低的陨石坑坑底（靛蓝色）与最高的山峰（白色）之间的高程差为 15 公里（10 英里）$^{[75]}$。“Ysolo 山”已被更名为 “Yamor 山”$^{[74]}$。} \label{fig_Ceres_9}
\end{figure}
“黎明号”探测器揭示，谷神星表面存在大量陨石坑，尽管大型陨石坑的数量比预期要少$^{[76]}$。基于当前小行星带形成模型的预测认为，谷神星应当有 10–15 个直径超过 400 公里（250 英里）的陨石坑$^{[76]}$。谷神星上已确认的最大陨石坑为 Kerwan 盆地，直径为 284 公里（176 英里）$^{[77]}$。最可能的原因是地壳发生黏性松弛，使得较大的撞击坑随着时间缓慢变平$^{[76][78]}$。

谷神星的北极区域显示出比赤道区域更多的陨石坑，尤其是东部赤道区域的撞击坑相对较少$^{[79]}$。直径在 20–100 公里（10–60 英里）之间的陨石坑总体数量分布与它们形成于晚期重轰炸时期一致，古老极地区域之外的陨石坑很可能在早期低温火山活动中被抹除$^{[79]}$。三个大型浅盆地（大平原），其边缘已退化，很可能是遭受侵蚀的陨石坑$^{[64]}$。其中最大的 Vendimia 大平原直径达 800 公里（500 英里）$^{[76]}$，同时也是谷神星上最大的单一地理特征$^{[80]}$。这三者之中有两个区域的铵浓度高于平均水平$^{[64]}$。

“黎明号”在谷神星表面观测到 4,423 块直径大于 105 米（344 英尺）的巨石。这些巨石很可能由撞击形成，并出现在陨石坑内或附近，但并非所有陨石坑都包含巨石。大型巨石在高纬度地区数量更多。谷神星上的巨石较为脆弱，并会由于热应力（在黎明和黄昏时，表面温度迅速变化）及陨石撞击而快速降解。它们的最大年龄估计约为 1.5 亿年，远短于灶神星上巨石的存续时间$^{[81]}$。

\textbf{构造特征}

尽管谷神星缺乏板块构造$^{[82]}$，其绝大多数表面特征要么与撞击有关，要么与低温火山活动有关$^{[83]}$，但其表面仍初步识别出一些潜在的构造特征，尤其是在其东半球。Samhain 坑链为谷神星表面尺度达数公里的线状裂隙，它们与撞击没有明显联系，反而更类似于坑链结构，暗示存在埋藏的正断层。此外，谷神星上若干陨石坑的浅层裂纹坑底与低温岩浆侵入相一致$^{[84]}$。

\textbf{低温火山活动}
\begin{figure}[ht]
\centering
\includegraphics[width=6cm]{./figures/998702449cec5764.png}
\caption{阿胡纳山（中央）在最陡的一侧估计有 5 公里（3 英里）高$^{[85]}$。} \label{fig_Ceres_11}
\end{figure}
\begin{figure}[ht]
\centering
\includegraphics[width=6cm]{./figures/7c3452ad060f1bc2.png}
\caption{谷神星上 Cerealia 亮斑与 Vinalia 亮斑的模拟视图} \label{fig_Ceres_12}
\end{figure}
谷神星有一座显著的山峰——阿胡纳山；该山似乎是一座低温火山，并且表面陨石坑很少，表明其最大年龄约为 2.4 亿年$^{[86]}$。其相对较高的重力场表明它密度较大，因此由岩石而非冰构成，而且其位置很可能是由于由盐水与硅酸盐颗粒组成的浆状物在地幔顶部发生底侵作用所致$^{[51]}$。它与 Kerwan 盆地大致呈对跖位置。形成 Kerwan 的撞击所产生的地震能量可能集中在谷神星的另一侧，破裂地壳外层，并触发高黏度低温岩浆（由含盐的泥状冰软化而成）向表面迁移$^{[87]}$。Kerwan 也显示出因撞击导致的地下冰融化而出现液态水作用的证据$^{[77]}$。

一项 2018 年的计算机模拟显示，谷神星上的低温火山一旦形成，会因黏性松弛而在数亿年时间内逐渐消退。研究团队在谷神星表面识别出 22 个强有力的候选特征，可能是经历了松弛作用的低温火山$^{[86][88]}$。Yamor 山是一座古老的、带有撞击陨石坑的山峰，尽管年龄更久远，但由于位于谷神星北极区域，低温阻止了地壳的黏性松弛，因此其外观与阿胡纳山相似$^{[83]}$。模型显示，在过去十亿年中，谷神星平均每 5,000 万年形成一座低温火山$^{[83]}$。火山喷发可能与古老的撞击盆地有关，但在谷神星上分布并不均匀$^{[83]}$。该模型表明，与阿胡纳山的发现相反，谷神星的低温火山必须由密度远低于谷神星地壳平均值的物质构成，否则无法出现已观察到的黏性松弛现象$^{[86]}$。

大量谷神星陨石坑意外地拥有中心坑，这可能与低温火山过程有关；其他陨石坑则有中央峰$^{[89]}$。“黎明号”探测器观察到数百个亮斑（faculae），其中最亮的一处位于直径 80 公里（50 英里）的 Occator 陨石坑中央$^{[90]}$。Occator 中心的亮斑被命名为 Cerealia 亮斑$^{[91]}$；其东侧的亮斑群被命名为 Vinalia 亮斑$^{[92]}$。Occator 有一个宽 9–10 公里的坑，部分被中央穹丘填充。该穹丘形成于亮斑之后，很可能是地下储层结冰导致的，类似于地球北极地区的冰丘（pingos）$^{[93][94]}$。在 Cerealia 上空周期性出现薄雾，这支持了亮斑由某种气体逸出或冰升华所形成的假说$^{[95]}$。2016 年 3 月，“黎明号”在 Oxo 陨石坑首次发现了谷神星表面存在水冰的确凿证据$^{[96]}$。

2015 年 12 月 9 日，NASA 科学家报告称，谷神星亮斑可能由一种来自蒸发盐水的盐类构成，包含六水合硫酸镁（MgSO(_4\cdot6)H(_2)O）；这些亮斑也与富含氨的粘土有关$^{[97]}$。2017 年的近红外光谱分析显示，这些亮斑与大量碳酸钠（Na(_2)CO(_3)）以及少量氯化铵（NH(_4)Cl）或碳酸氢铵（NH(_4)HCO(_3)）相一致$^{[98][99]}$。这些物质被认为起源于到达地表的盐水结晶$^{[100]}$。2020 年 8 月，NASA 确认谷神星是一个富含水的天体，具有深层盐水储库，这些盐水在数百处渗出地表$^{[101]}$，形成“亮斑”，包括出现在 Occator 陨石坑中的那些$^{[102]}$。
\subsubsection{内部结构}
\begin{figure}[ht]
\centering
\includegraphics[width=6cm]{./figures/8a05aad22b023501.png}
\caption{谷神星内部结构的三层模型：外层厚壳（冰、盐类、水合矿物）,富含盐分的液体（盐水）与岩石,“地幔”（水合岩石）} \label{fig_Ceres_13}
\end{figure}
谷神星的活跃地质活动由冰和盐水驱动。从岩石中淋滤出的水估计含有约 5\% 的盐分。总体而言，谷神星按体积计约有 50\% 为水（相比之下地球为 0.1\%），按质量计约有 73\% 为岩石$^{[14]}$。

谷神星最大的一些陨石坑深达数公里，这与富含冰的浅层地下结构不相符。其表面能够保存直径接近 300 公里（200 英里）的陨石坑这一事实表明，谷神星最外层约比水冰强 1000 倍。这与一种由硅酸盐、水合盐和甲烷水合物混合物构成的模型一致，其中水冰体积比例不超过 30\%$^{[64][103]}$。

“黎明号”探测器的重力测量提出了谷神星内部结构的三种互相竞争的模型$^{[14]}$。在三层模型中，谷神星被认为由一个约 40 公里（25 英里）厚的外层地壳组成，其成分包括冰、盐类和水合矿物；内部为泥状的“地幔”，由水合岩石（例如粘土）构成，两者之间隔着一层约 60 公里（37 英里）厚的泥状盐水与岩石混合层$^{[104]}$。目前无法确定谷神星深部内部是否含有液态物质，或是否存在一个富含金属的高密度核心$^{[104]}$，但其较低的中心密度表明其内部可能保留约 10\% 的孔隙度$^{[14]}$。一项研究估计，核层与地幔/地壳层的密度分别为 2.46–2.90 与 1.68–1.95 g/cm³，而地幔与地壳合计厚度为 70–190 公里（40–120 英里）。预计核心只会发生部分脱水（冰的排出），但地幔相对于水冰的高密度反映其富含硅酸盐和盐类$^{[9]}$。也就是说，核心（若存在）、地幔和地壳均由岩石和冰组成，只是比例不同。

谷神星的矿物组成只能（间接地）确定其外层约 100 公里（60 英里）的区域。其固体外壳厚约 40 公里（25 英里），为冰、盐类和水合矿物的混合物。其下可能存在一薄层含有少量盐水，该结构延伸至至少 100 公里（60 英里）的探测深度以下，再往下被认为是主要由水合岩石（如粘土）构成的地幔$^{[104]}$。

在一种两层模型中，谷神星由一个球粒状物质组成的核心和一个由冰与微米级固体微粒（即“泥”）构成的地幔组成。表面冰的升华将留下约 20 米厚的水合微粒沉积层。区分程度的范围与数据相符：从一个大型核心（直径 360 公里〔220 英里〕，其中 75\% 为球粒、25\% 为微粒，地幔成分为 75\% 冰、25\% 微粒），到一个小型核心（直径 85 公里〔55 英里〕，几乎完全由微粒组成，地幔为 30\% 冰、70\% 微粒）。若核心较大，则核心—地幔边界温度应足够高，可形成盐水囊；若核心较小，则在 110 公里（68 英里）以下地幔应保持为液态。在后者情况下，液体储层中 2\% 的冻结就足以产生压力迫使部分液体上升至表面，造成低温火山活动$^{[105]}$。

另一种两层模型认为谷神星发生了部分分异，即形成一个富含挥发物的地壳与一个较致密的水合硅酸盐地幔。根据可能撞击过谷神星的陨石类型，可以计算一系列可能的地壳与地幔密度。若假定为 CI 类陨石（密度 2.46 g/cm³），则地壳厚度约 70 公里（40 英里），密度为 1.68 g/cm³；若为 CM 类陨石（密度 2.9 g/cm³），则地壳厚度约 190 公里（120 英里），密度为 1.9 g/cm³。拟合最佳的模型得出地壳厚度约 40 公里（25 英里），密度约 1.25 g/cm³，而地幔/核心的密度约为 2.4 g/cm³$^{[64]}$。
\subsection{外逸层}
2017 年，“黎明号”确认谷神星存在短暂的水蒸气大气$^{[106]}$。早在 2014 年初就已出现大气迹象，当时赫歇尔空间天文台在谷神星上探测到局部的中纬度水蒸气来源，其范围不超过 60 公里（40 英里），每秒释放约 (10^{26}) 个（3 千克）水分子$^{[107][108]}$。两个潜在来源区域——Piazzi 区（东经 123°，北纬 21°）与 A 区（东经 231°，北纬 23°）——在 W. M. Keck 天文台的近红外观测中显示为暗区（A 区具有明亮中心）。可能的蒸汽释放机制包括约 0.6 平方公里（0.2 平方英里）裸露表面冰的升华、由放射性内部热驱动的低温火山喷发$^{[107]}$，或由于上层冰变厚导致地下海洋受压$^{[111]}$。2015 年，David C. Jewitt 将谷神星列入活跃小行星名单$^{[112]}$。在距太阳小于 5 AU 的位置，表面冰处于不稳定状态$^{[113]}$，因此直接暴露在太阳辐射下预期会升华。太阳耀斑和日冕物质抛射所产生的质子会溅射表面裸露冰，使得水蒸气探测与太阳活动呈正相关$^{[114]}$。水冰可以从谷神星深部迁移至表面，但会在短时间内逃逸。当谷神星远离太阳时，表面升华预期会降低，而内部驱动的释放则不受轨道位置影响。先前有限的数据暗示谷神星升华类似彗星$^{[107]}$，但“黎明号”的证据表明地质活动至少部分负责这一现象$^{[115]}$。

利用“黎明号”的伽马射线与中子探测器（GRaND）进行的研究显示，谷神星能加速来自太阳风的电子；最被接受的假说认为，这些电子通过太阳风与稀薄水蒸气外逸层的碰撞被加速$^{[116][117]}$。这种弓形激波也可能由一个短暂的磁场解释，但这种可能性较小，因为谷神星内部被认为没有足够的电导率$^{[117]}$。谷神星的稀薄外逸层通过撞击暴露的冰斑、冰通过多孔冰壳的扩散以及太阳活动期间质子溅射不断得到补充$^{[118][119][120]}$。蒸汽扩散速率随颗粒大小变化$^{[121]}$，并受一个由约 1 微米粒子聚集而成的全球尘埃覆盖层的强烈影响$^{[122]}$。仅靠升华实现的外逸层补给非常微小，目前的气体释放率仅为 0.003 千克/秒$^{[119]}$。有关外逸层存在的多种模型已被提出，包括弹道轨迹模型、DSMC 模型与极地帽数值模型$^{[123][120][124]}$。结果显示，弹道轨迹模型得到的水外逸层半衰期为 7 小时；DSMC 模型在光学稀薄大气下得到了持续数十天且释放率为 6 千克/秒的情形；极地帽模型显示，由外逸层输送的水会形成季节性极地冰盖。外逸层中水分子的运动主要受弹跳式运动与表面相互作用的支配，但对其与行星风化层直接相互作用的了解较少$^{[119]}$。
\subsection{起源与演化}
谷神星是一颗幸存的原行星，形成于 45.6 亿年前；与智神星（Pallas）与灶神星一起，是太阳系内层中仅存的三颗原行星之一$^{[125]}$，其余要么在碰撞中合并形成类地行星，要么在撞击中被破碎$^{[126]}$，或被木星抛出轨道$^{[127]}$。尽管谷神星当前的位置位于小行星带，但其成分并不符合在该区域形成的特征。相反，谷神星似乎形成于木星与土星轨道之间，并在木星向外迁移时被偏转进入小行星带$^{[14]}$。在 Occator 陨石坑中发现的铵盐支持这一源自外太阳系的假说，因为氨在外太阳系丰富得多$^{[128]}$。

谷神星早期地质演化依赖于其形成期间与形成后的热源：来自星子聚积的撞击能量以及放射性核素（可能包括已灭绝的短寿命核素，如 (^{26})Al）的衰变。这些热源可能足以使谷神星分化形成岩质行星核与冰质地幔，甚至形成液态海洋$^{[64]}$，这一过程在其形成后不久就可能发生$^{[66]}$。这一海洋在冻结时应在表面之下留下冰层。“黎明号”没有发现这类证据，表明谷神星原始地壳至少部分被后期撞击破坏，使得冰与盐类以及富含硅酸盐的古代海底物质深度混合$^{[64]}$。

谷神星拥有令人惊讶的少量大型陨石坑，表明黏性松弛与低温火山活动已抹去了较古老的地质特征$^{[129]}$。粘土与碳酸盐的存在需要高于 50°C 的化学反应，符合热液活动条件$^{[51]}$。

谷神星的地质活动随时间显著减弱，其表面主要由撞击坑主导；然而，“黎明号”的证据显示，内部过程仍持续在相当程度上塑造谷神星表面$^{[130]}$，这与此前预测谷神星因其较小体积而在早期即终止内部地质活动的观点相反$^{[131]}$。
\subsection{可居性}
\begin{figure}[ht]
\centering
\includegraphics[width=6cm]{./figures/7721eaf63a85b575.png}
\caption{} \label{fig_Ceres_14}
\end{figure}
 